\section*{第 \ref{ch:b2:affine} 章——仿射空间}

\begin{solution}{exo:b2:affine:1}
$\{x + y + z = 1\}$：仿射平面，方向为向量平面 $\{x
+ y + z = 0\}$，维数为 $2$。添加 $x - z = 3$ 后得仿射直线
（两个独立方程），方向为 $\{x + y + z = 0,\ x = z\}
= \operatorname{Vect}\bigl((1, -2, 1)\bigr)$，维数为 $1$。$\{x^2
+ y^2 = 1\}$：柱面——对\omterm{def:b2:affine:barycenter}{重心}运算不封闭（
$(1,0,0)$ 与 $(-1,0,0)$ 的中点为原点，不在
柱面上）：故非仿射。相容方程组 $MX = B$ 的解集为仿射
子空间 $X_0 + \ker M$，维数为 $\dim\ker M$，如
\cref{def:b2:affine:subspace} 所回顾。
\end{solution}

\begin{solution}{exo:b2:affine:2}
$C = \operatorname{bar}(A, 1-t; B, t)$ 表示 $\vect{AC} =
t\,\vect{AB}$：$t$ 存在当且仅当
$\vect{AC}$ 与 $\vect{AB} \neq 0$ 线性相关，即三点共线。位置如下：$t
= \frac12$：中点；$t = 2$：位于 $B$ 之外，与 $B$ 的距离等于 AB 的长
（$\vect{AC} = 2\vect{AB}$）；$t = -1$：$B$ 关于
$A$ 的对称点。
\end{solution}

\begin{solution}{exo:b2:affine:3}
$M$ 为到 $\operatorname{Vect}(1,1)$ 上的投影矩阵，投影沿
$(1,-1)$ 方向（验证 $M^2 = M$）。$f$ 的像：$\{MX + C\} = C +
\operatorname{im} M$：过 $(1,0)$，方向为
$(1,1)$ 的仿射直线。不动点：$X = MX + C$，即 $(I - M)X = C$；但 $C =
(1, 0)^{\mathsf T}$，且 $\operatorname{im}(I - M) =
\operatorname{Vect}(1,-1)$；$(1,0)$ 在其中吗？$(1, 0) =
\alpha(1,-1)$ 推出 $\alpha = 1$ 且 $0 = -1$，不可能。故无
不动点。又
\[
f(f(X)) = M(MX + C) + C = MX + MC + C = f(X) + MC,
\qquad MC = \tfrac12(1,1)^{\mathsf T} :
\]
$f\circ f$ 相当于先作 $f$，再沿像直线作平移，
故 $f$ 为“滑动投影”：先投影到直线上，
再与一个滑动复合。
\end{solution}

\begin{solution}{exo:b2:affine:4}
$I = \operatorname{bar}(B, 2; C, 1)$（由于 $\vect{BI} =
\frac13\vect{BC}$ 使 $I$ 靠近 $B$，$2$ 为 $B$ 上的权，
$1$ 为 $C$ 上的权——验证：$\vect{BI} = \frac{1}{3}\vect{BC}$）。
同理 $J = \operatorname{bar}(C, 2; A, 1)$，$K =
\operatorname{bar}(A, 2; B, 1)$。将三个带权点组相加，
则 $(I, 1; J, 1; K, 1)$ 的\omterm{def:b2:affine:barycenter}{重心}（每个的总权为
$3$，故将 $I$ 换为其带权点组，依此类推）为
\[
\operatorname{bar}\bigl(A, 1 + 2;\ B, 2 + 1;\ C, 1 + 2\bigr)
= \operatorname{bar}(A, 1; B, 1; C, 1) = G ,
\]
即 $ABC$ 的\omterm{def:b2:affine:barycenter}{重心}：三角形 $IJK$ 与之有相同\omterm{def:b2:affine:barycenter}{重心}——
这是可以预料的，因为构造关于 $A, B, C$ \omterm{def:b2:structures:sn}{轮换}对称，
而\omterm{def:b2:affine:barycenter}{重心}是此\omterm{def:b2:structures:sn}{轮换}对称下的唯一不动点，
这由权的对称性决定。
\end{solution}

\begin{solution}{exo:b2:affine:5}
设 $g(u) = f(O + u) - f(O)$（在 $\R^n$ 中，以
$O$ 为基点向量化），$g(0) = 0$。

\emph{可加性}：$\frac{(O + u) + (O + v)}{2} = O + \frac{u +
v}{2}$，由保中点性得 $g\bigl(\frac{u+v}{2}\bigr)
= \frac{g(u) + g(v)}{2}$；取 $v = 0$，得 $g(u/2) = g(u)/2$；
综合可得 $g(u + v) = 2g\bigl(\frac{u+v}{2}\bigr) = g(u) + g(v)$。

\emph{$\Q$-齐次性}：由可加性得 $g(nu) = ng(u)$（$n \in
\N$，由归纳法），再得 $g(-u) = -g(u)$（相加即得），再得 $g(\frac pq u) =
\frac pq g(u)$（作用 $q$ 次，并利用数乘的单射性）。

\emph{$\R$-齐次性}：对 $t \in \R$，取有理数列 $t_n \to
t$：$g(t_nu) = t_ng(u)$，由 $g$（继承自 $f$）的\omterm{def:b2:metric:continuity}{连续}性
取极限即得 $g(tu) = tg(u)$。从而 $g$ 是线性的，且 $f =
f(O) + g$，故为\omterm{def:b2:affine:subspace}{仿射映射}。
\end{solution}

\begin{solution}{exo:b2:affine:6}
线性部分 $\varphi(x,y) = (y, x)$：关于对角线的反射
$y = x$（正交，\omterm{def:b2:linalg:det}{行列式}为 $-1$）。不动点：$(x, y) = (y
+ 1, x - 1)$ 等价于单个方程 $y = x - 1$（两个分量等价）：
直线 $y = x - 1$ 上每个点都是不动点。
因此 $f$ 逐点固定该直线，即 $f$
是该轴上的\emph{反射}（一个等距变换，不动点构成直线，
线性部分为反射）。一致地，$f \circ
f(x,y) = f(y+1, x-1) = (x - 1 + 1, y + 1 - 1) = (x, y)$，即
对合（对合性正是反射所应有的）。
\end{solution}

\begin{solution}{exo:b2:affine:7}
$n + 1$ 个向量 $\vect{A_1A_i}$（$i \geq 2$）在
$n$ 维空间中线性相关：存在不全为零的 $\mu_i$，使 $\sum_{i\geq2}
\mu_i\vect{A_1A_i} = 0$；令 $\mu_1 = -\sum_{i \geq 2}\mu_i$，则
$\sum_{i}\mu_i = 0$，且 $\sum_i \mu_i\,\vect{OA_i} = 0$ 对每个
$O$ 成立，其中 $\mu_i$ 不全为零。将下标分成 $P = \{i : \mu_i >
0\}$，$N = \{i : \mu_i < 0\}$，两者皆非空（$\mu_i$ 之和为
零且不全为零）。记 $s = \sum_{i\in P}\mu_i =
-\sum_{i \in N}\mu_i > 0$，有
\[
\operatorname{bar}\bigl(A_i, \tfrac{\mu_i}{s}\bigr)_{i \in P}
= \operatorname{bar}\bigl(A_i, \tfrac{-\mu_i}{s}\bigr)_{i \in N},
\]
（两边都等于点 $X$，其中 $\vect{OX} = \frac1s\sum_{i\in
P}\mu_i\vect{OA_i}$，由该关系式）：这是两个
\omterm{def:b2:affine:convex}{凸包}的公共点，且两组下标不相交。
\end{solution}

\begin{solution}{exo:b2:affine:8}
\emph{像即不动点}：对 $Y = f(X)$，有 $f(Y) = f(f(X)) =
f(X) = Y$，故每个像点都是不动点；反之，不动点都是像。
于是 $\mathcal{F} = \operatorname{im} f =
\operatorname{Fix}(f)$ 非空，且它是仿射子空间
（\omterm{def:b2:affine:subspace}{仿射映射}的像），其方向为 $\operatorname{im}\vec f$。

\emph{投影结构：}$\vec f$ 是幂等的（$\vec{f\circ
f} = \vec f^{\,2} = \vec f$），故 $E = \operatorname{im}\vec f
\oplus \ker \vec f$（见\cref{ex:b2:reduction:projections}）。对任意
点 $X$，考虑向量 $\vect{f(X)\,X}$；用 $\vec f$ 作用其上：
\[
\vec f\bigl(\vect{f(X)\,X}\bigr) = \vect{f(f(X))\,f(X)} = 0
\qquad (f \circ f = f),
\]
故 $\vect{f(X)\,X} \in \ker\vec f$。于是 $X = f(X) +
\vect{f(X)X}$ 表明 $X$ 可视为 $\mathcal{F}$ 中的点
经 $\ker\vec f$ 中某向量平移所得，故 $f$ 恰是
到 $\mathcal{F}$ 上沿 $\ker\vec f$ 的投影。反之，
此类投影显然满足 $f \circ f = f$。
\end{solution}

\begin{solution}{exo:b2:affine:9}
设 $M = \operatorname{bar}(B, 2;\ C, 3)$，总权为 $5$。
由结合性得 $G = \operatorname{bar}(A, 1;\ M, 5)$，故
$\vect{AG} = \frac56\,\vect{AM}$，即 $G$ 位于线段
$\intcc AM$ 上，距 $A$ 为其全长的六分之五处。由 $A \notin
(BC)$，则直线 $(AG) = (AM)$ 与 $(BC)$ 唯一相交于点
$M$，且 $\vect{BM} = \frac35\,\vect{BC}$。同理，取 $N =
\operatorname{bar}(C, 3;\ A, 1)$（总权为 $4$，$\vect{CN}
= \frac14\,\vect{CA}$），由结合性得 $G =
\operatorname{bar}(B, 2;\ N, 4)$，直线 $(BG)$ 与 $(CA)$ 相交于
$N$，且 $\vect{BG} = \frac46\,\vect{BN} =
\frac23\,\vect{BN}$。
\end{solution}

\begin{solution}{exo:b2:affine:10}
取原点 $O$，把各点写成向量：
$h_{\Omega, \lambda}(x) = \omega + \lambda(x - \omega)$，其中
$\omega = \vect{O\Omega}$。复合映射 $g = h_{\Omega',
\mu} \circ h_{\Omega, \lambda}$ 为\omterm{def:b2:affine:subspace}{仿射映射}，其线性部分为
$\mu\lambda\,\mathrm{id}$。若 $\lambda\mu \neq 1$，则 $1 \notin
\operatorname{Sp}(\lambda\mu\,\mathrm{id})$，故
由\cref{prop:b2:affine:fixedpoint}可知唯一不动点
$\Omega''$，在该点向量化，得 $g =
\lambda\mu\,\mathrm{id}$，即位似变换 $h_{\Omega'',
\lambda\mu}$。若 $\lambda\mu = 1$，则线性部分为
恒等映射，故 $g$ 为平移；展开得
\[
g(x) = \omega' + \mu\bigl(\omega + \lambda(x - \omega) -
\omega'\bigr) = x + (1 - \mu)\,\omega' + \mu(1 -
\lambda)\,\omega .
\]
当 $\lambda = \mu = -1$（点反射）时，该向量为
$2\omega' - 2\omega = 2\,\vect{\Omega\Omega'}$，即
依次关于 $\Omega$、再关于
$\Omega'$ 的点反射的复合，为平移 $2\,\vect{\Omega\Omega'}$。
\end{solution}

\begin{solution}{exo:b2:affine:11}
$\vect{A'B} = \alpha\,\vect{A'C}$ 恰等价于 $1\cdot
\vect{A'B} - \alpha\,\vect{A'C} = 0$，即 $A' =
\operatorname{bar}(B, 1;\ C, -\alpha)$（总权为 $1 -
\alpha \neq 0$，因 $B \neq C$）；同理 $B' =
\operatorname{bar}(C, 1;\ A, -\beta)$，以及 $C' =
\operatorname{bar}(A, 1;\ B, -\gamma)$。

\emph{共线判据。}取各点 $P$ 的
规范化\omterm{def:b2:affine:barycenter}{重心}坐标行向量 $p = (p_A, p_B, p_C)$，$p_A + p_B +
p_C = 1$，相对于 $(A, B, C)$。若 $\sum_i c_i p_i = 0$
其中 $(c_1, c_2, c_3) \neq 0$，对三点 $P_1, P_2,
P_3$，则对各分量求和得 $\sum c_i = 0$，且
$\sum_i c_i \vect{OP_i} = \sum_j \bigl(\sum_i
c_ip_{ij}\bigr)\vect{OV_j} = 0$：各 $P_i$ 仿射相关，
即共线。反之，仿射相关式
$(t_i)$ 给出 $w = \sum t_ip_i$，其分量之和为 $0$，
且 $\sum_j w_j\vect{OV_j} = 0$；以 $A$ 为基点展开，$w_B
\vect{AB} + w_C\vect{AC} = 0$，故 $w = 0$，由
$(A, B, C)$ 仿射无关，故各行线性相关。
因此共线等价于 $3 \times 3$ \omterm{def:b2:linalg:det}{行列式}为零，
而将各行乘以非零因子 $1 - \alpha$、$1 -
\beta$、$1 - \gamma$ 不影响结论：
\[
\det\begin{pmatrix} 0 & 1 & -\alpha\\ -\beta & 0 & 1\\ 1 &
-\gamma & 0\end{pmatrix}
= 1 - \alpha\beta\gamma .
\]
因此 $A', B', C'$ 共线，当且仅当 $\alpha\beta\gamma = 1$：
此即梅涅劳斯定理。
\end{solution}

\begin{solution}{exo:b2:affine:12}
设 $\Delta = \{\lambda \in \R^{n+1} : \lambda_i \geq 0,\
\sum\lambda_i = 1\}$，它在 $\R^{n+1}$ 中闭且有界，故
为紧集，且 $K^{n+1}$ 作为有限积也是紧的。映射
\[
\Phi \colon \Delta \times K^{n+1} \to \R^n, \qquad
\Phi(\lambda, x_0, \dots, x_n) = \sum_{i=0}^n \lambda_i x_i
\]
\omterm{def:b2:metric:continuity}{连续}，而\cref{thm:b2:affine:caratheodory}断言
恰好 $\operatorname{conv}(K) = \Phi(\Delta \times
K^{n+1})$：紧集的\omterm{def:b2:metric:continuity}{连续}像
（\cref{thm:b2:metric:compactprops}），故为紧集。

对闭集，取 $S = (\R \times \{0\}) \cup \{(0,
1)\}$，在 $\R^2$ 中是闭集。若凸组合将权重
$t$ 置于 $(0,1)$，$1 - t$ 置于轴上的点，则其第二坐标为
$t$，故
\[
\operatorname{conv}(S) = \bigl(\R \times \intco01\bigr) \cup
\{(0,1)\}
\]
（对 $0 \leq t < 1$，$(x, t) = t\,(0,1) + (1-t)\,\bigl(\tfrac
x{1-t}, 0\bigr)$）。点 $(1, 1) = \lim_{t \to 1}(1, t)$
是附着点，但不在\omterm{def:b2:affine:convex}{凸包}中，故不闭。
\end{solution}

\begin{solution}{pb:b2:affine:1}
\textbf{1.} 以 $A_j$ 为基点：对 $i \neq j$，$\vect{A_jA_i} =
\vect{A_0A_i} - \vect{A_0A_j}$。若 $\sum_{i \neq j}
c_i\vect{A_jA_i} = 0$，展开得 $\sum_{i \notin \{0,
j\}} c_i\,\vect{A_0A_i} - \bigl(\sum_{i\neq j}
c_i\bigr)\vect{A_0A_j} = 0$；由
$\vect{A_0A_i}$ 的线性无关性，得 $c_i = 0$（$i \notin \{0, j\}$），再
$c_0 = 0$，故在 $A_j$ 处仿射无关。关于等价性：两个
权族 $(\lambda_i)$、$(\lambda_i')$ 之和均为 $1$，
有相同\omterm{def:b2:affine:barycenter}{重心}，取 $\nu = \lambda -
\lambda'$，则 $\sum\nu_i = 0$，且（取原点 $A_0$）$\sum_{i \geq
1}\nu_i\,\vect{A_0A_i} = 0$，由线性无关得 $\nu = 0$。
反之，若存在非平凡关系 $\sum_{i\geq1}\mu_i
\vect{A_0A_i} = 0$，补取 $\mu_0 = -\sum_{i\geq1}
\mu_i$，可对任一权族加上 $t(\mu_i)$ 而不改变
\omterm{def:b2:affine:barycenter}{重心}，故\omterm{def:b2:affine:barycenter}{重心}表示不唯一。

\textbf{2.} $(\vect{AB}, \vect{AC})$ 构成 $\R^2$ 的一组基：
写 $\vect{AM} = \beta\,\vect{AB} + \gamma\,\vect{AC}$
（唯一），并令 $\alpha = 1 - \beta - \gamma$；\omterm{def:b2:affine:barycenter}{重心}
条件在原点 $A$ 处恰为 $\vect{AM} =
\beta\,\vect{AB} + \gamma\,\vect{AC}$。唯一性即
第 1 问。

\textbf{3.} 由 $\alpha\vect{MA} + \beta\vect{MB} +
\gamma\vect{MC} = 0$ 及夏勒关系式，得 $\vect{MA} =
-\beta\,\vect{AB} - \gamma\,\vect{AC}$。记 $D =
\det(\vect{AB}, \vect{AC})$：
\begin{align*}
\det(\vect{MB}, \vect{MC})
&= \det(\vect{MA} + \vect{AB},\ \vect{MA} + \vect{AC})\\
&= \det(\vect{MA}, \vect{AC}) + \det(\vect{AB}, \vect{MA}) +
D = -\beta D - \gamma D + D = \alpha D,
\end{align*}
利用双线性及 $\det(\vect{MA}, \vect{AC}) = -\beta D$，
$\det(\vect{AB}, \vect{MA}) = -\gamma D$。另外两式
可由同样的计算得出，只需将角色
依\omterm{def:b2:structures:sn}{轮换}次序置换即可。

\textbf{4.} 由定义，$\operatorname{conv}\{A, B, C\}$ 是
非负权\omterm{def:b2:affine:barycenter}{重心}全体构成的集合；将权规范化，使其和为
$1$，并利用唯一性（第 2 问），
$M \in \operatorname{conv}\{A,B,C\}$ 当且仅当 $\alpha, \beta,
\gamma \geq 0$。每个坐标都是 $M$ 的仿射函数
（第 3 问：$2\times2$ \omterm{def:b2:linalg:det}{行列式}，其一列
关于 $M$ 仿射），故每个开符号条件都定义一个开
半平面。模式 $(-,-,-)$ 与 $\alpha + \beta
+ \gamma = 1$ 矛盾；其余七种模式均可实现：
取符号相符且至少有一个
$+$ 分量的三元组，将其缩放，使（正）和为 $1$——例如
$(-1, 1, 1)$、$(3, -1, -1)$、$(\frac13, \frac13, \frac13)$、
以及各种置换。

\textbf{5.} \omterm{def:b2:affine:subspace}{仿射映射}保持\omterm{def:b2:affine:barycenter}{重心}
（见\cref{def:b2:affine:subspace}），故 $u(M) = \alpha u(A) +
\beta u(B) + \gamma u(C)$。记 $u(x, y) = ax + by + c$
其中 $(a, b) \neq (0,0)$：$\{u = c'\}$ 为一条直线，且每条
直线 $ax + by = c'$ 都是这样的水平集。若 $u(M), u(N) \geq
c$，$t \in \intcc01$，则 $u\bigl(\operatorname{bar}(M,
1-t; N, t)\bigr) = (1-t)u(M) + tu(N) \geq c$：半平面是
凸的。

\textbf{6.} 条件 $\sum\mu_i = 0$，$3\mu_2 + \mu_4 =
0$，$3\mu_3 + \mu_4 = 0$ 给出（取 $\mu_4 = 3$）的
相关关系 $(\mu_1, \mu_2, \mu_3, \mu_4) = (-1, -1, -1, 3)$。
符号划分为 $\{A_4\} \mid \{A_1, A_2, A_3\}$，并将
两侧各除以 $3$，以规范化：
\[
A_4 = \operatorname{bar}\bigl(A_1, \tfrac13;\ A_2, \tfrac13;\
A_3, \tfrac13\bigr) = (1,1),
\]
这是三角形的\omterm{def:b2:affine:barycenter}{重心}：拉东点为 $A_4$
本身，确实位于三角形 $A_1A_2A_3$ 内部。

\textbf{7.} 线性映射 $\Phi \colon \R^{n+2} \to \R
\times \R^n$，$\mu \mapsto (\sum\mu_i,\
\sum\mu_i\vect{OA_i})$，其秩 $\leq n + 1$，故 $\dim\ker
\Phi \geq 1$。若两个线性无关的相关关系 $\mu, \mu'$
存在，则取适当组合 $\nu = \mu'_{n+2}\mu -
\mu_{n+2}\mu'$（若两末项系数均为零，则取
$\mu$ 本身）即为一个\emph{非零}相关关系，且 $\nu_{n+2}
= 0$；限制于 $A_1, \dots, A_{n+1}$ 上，并以
$A_1$ 为基点，则存在 $\nu_i \neq 0$，其中 $i \geq 2$（单个非零权
不可能和为零），由此得非平凡关系
$\sum_{i\geq2}\nu_i\vect{A_1A_i} = 0$，这 $n+1$ 个点
将仿射相关，与一般位置矛盾。故 $\dim\ker
\Phi = 1$。同样的限制论证表明，非零
相关关系的系数均不为零。

\textbf{8.} 设 $\mu \neq 0$ 为相关关系，$P = \{i :
\mu_i > 0\}$ 与 $N = \{i : \mu_i < 0\}$：二者均非空
（$\sum\mu_i = 0$，$\mu \neq 0$），且穷尽（无零
系数）。拉东构造（见\cref{exo:b2:affine:7}）
恰由此划分得到两\omterm{def:b2:affine:convex}{凸包}的公共点。
由于相关关系在非零数乘意义下唯一
（第 7 问），无序对 $\{P, N\}$——从而
拉东划分——是唯一的。

\textbf{9.} 各块非空，故类型为 $(1,3)$ 或
$(2,2)$。类型 $(1,3)$，块 $\{j\}$：拉东点位于
$\operatorname{conv}\{A_j\} = \{A_j\}$ 中，故 $A_j \in
\operatorname{conv}$，即属于其余三点的\omterm{def:b2:affine:convex}{凸包}；它不能落在一条
边上（否则三点共线，与一般位置矛盾），
故 $A_j$ 位于三角形内部。类型
$(2,2)$，块 $\{i,j\} \mid \{k,l\}$：拉东点 $z$
位于 $\intcc{A_i}{A_j} \cap \intcc{A_k}{A_l}$ 上，且 $z$
不是端点（否则会使三点共线），两线段
在内部一点相交。此外，任一点都不落在
其余点的\omterm{def:b2:affine:convex}{凸包}内：此类包含 $A_l =
\operatorname{bar}(A_i, \lambda_i)_{i \neq l}$，其中
$\lambda_i \geq 0$，它是符号模式为
$(+,+,+,-)$ 的仿射相关关系，由唯一性（第 8 问）将使
划分为 $(1,3)$。因此在 $(2,2)$ 情形下，四点
处于凸位置，且相交线段即为
对角线。

\textbf{10.} 由方程 $\mu_2 + \mu_3 = 0$，$\mu_3 +
\mu_4 = 0$，$\sum\mu_i = 0$ 得到依赖关系 $(1, -1, 1,
-1)$：划分 $\{(0,0), (1,1)\} \mid \{(1,0), (0,1)\}$，且
\[
\operatorname{bar}\bigl((0,0), \tfrac12;\ (1,1),
\tfrac12\bigr) = \bigl(\tfrac12, \tfrac12\bigr) =
\operatorname{bar}\bigl((1,0), \tfrac12;\ (0,1),
\tfrac12\bigr) :
\]
拉东点是正方形的中心，两条
对角线在此相交——属 $(2,2)$ 型，与图中所示一致。

\textbf{11.} 对 $x_1, \dots, x_4$ 应用拉东定理，得到两块
$I \mid J$ 及一点 $z \in \operatorname{conv}\{x_i : i
\in I\} \cap \operatorname{conv}\{x_j : j \in J\}$。固定 $k
\in \{1, \dots, 4\}$，不妨设 $k \in I$。对每个 $j \in J$
满足 $j \neq k$，故 $x_j \in C_k$，由选取条件 $x_j \in
\bigcap_{l \neq j}C_l$。由于 $C_k$ 是\omterm{def:b2:affine:convex}{凸集}，$z \in
\operatorname{conv}\{x_j : j \in J\} \subseteq C_k$。因 $k$
任取，故 $z \in C_1 \cap C_2 \cap C_3 \cap C_4$。

\textbf{12.} 对 $m$ 用归纳法。当 $m = 3$ 时，假设即为
结论；$m = 4$ 即第 11 问。设 $m \geq 4$，
假设结论对 $m$ 个集合成立，并取 $C_1, \dots,
C_{m+1}$ 满足三交性质。令 $C_m' =
C_m \cap C_{m+1}$，为\omterm{def:b2:affine:convex}{凸集}。族 $C_1, \dots, C_{m-1},
C_m'$ 共 $m$ 个集合；不含 $C_m'$ 的三元组由归纳假设相交，
且三元组 $\{C_i, C_j, C_m'\}$ 的交为
$C_i \cap C_j \cap C_m \cap C_{m+1}$，非空，
由第 11 问用于 $C_i, C_j, C_m, C_{m+1}$（这四者中
任三者相交，由假设）。归纳假设现于新族给出
一个公共点，即所有 $m+1$
个集合的公共点。

\textbf{13.} (a) 非退化三角形的闭边 $\intcc PQ$、$\intcc QR$，
$\intcc RP$ 中任两条共享一个顶点，
但三者若有公共点，则该点须位于 $\intcc
PQ \cap \intcc RP = \{P\}$ 中，且在 $\intcc QR$ 中，
$P$ 不在其中，矛盾。(b) 任取三个集合 $S_i = \{x_1, \dots, x_4\}
\setminus \{x_i\}$ 各省去四点中的三点，余下
恰一个公共点；全体的交集不含每一点。
$S_i$ 为有限集而非\omterm{def:b2:affine:convex}{凸集}：凸性
不可或缺。(c) 有限个 $H_k = \intco{k}{+\infty} \times
\R$ 的交为 $\intco{k_{\max}}{+\infty} \times \R \neq
\emptyset$，然而不存在点使 $x \geq k$ 对所有 $k \in \N$ 皆成立：
对无限族而言，紧性不可或缺。

\textbf{14.} 假设 $\bigcap_{i \in I}K_i = \emptyset$，
并固定 $i_0$。每个 $x \in K_{i_0}$ 必不属于某个 $K_i$，故
$K_{i_0} \subseteq \bigcup_{i \in I}(\R^2 \setminus K_i)$，是一个
由\omterm{def:b2:metric:topology}{开集}构成的覆盖（$K_i$ 紧，从而为闭集）。由
博雷尔--勒贝格（\cref{thm:b2:metric:borellebesgue}）
知有限个即已足够：$K_{i_0} \cap K_{i_1} \cap \dots \cap K_{i_N} =
\emptyset$。但此有限族中任意三个\omterm{def:b2:affine:convex}{凸集}都相交，
故由第 12 问，该交集非空，
矛盾。

\textbf{15.} 令 $D_p = \overline D(p, r)$（$p \in S$）：
紧\omterm{def:b2:affine:convex}{凸集}。对 $p, q, s \in S$，由假设
知存在闭圆盘 $\overline D(z, r)$ 包含 $p, q, s$；
于是 $\norm{\vect{zp}}, \norm{\vect{zq}}, \norm{\vect{zs}}
\leq r$，即 $z \in D_p \cap D_q \cap D_s$。由赫利定理
（第 12 问；此族有限），存在 $c \in
\bigcap_{p \in S}D_p$：即对每个 $p \in S$ 满足
$\norm{\vect{cp}} \leq r$，故 $S \subseteq \overline D(c,
r)$。

\textbf{16.} 取 $c = x_{\lceil n/2\rceil}$。包含 $c$ 的闭
半直线为 $\intoc{-\infty}{t}$，其中 $t \geq
c$，或 $\intco{t}{+\infty}$，其中 $t \leq c$。第一者
包含 $x_1, \dots, x_{\lceil n/2\rceil}$，至少有
$\lceil n/2\rceil \geq n/2$ 个点。第二者包含
$x_{\lceil n/2\rceil}, \dots, x_n$，恰有 $n - \lceil
n/2\rceil + 1 = \floor{n/2} + 1 > n/2$ 个点。

\textbf{17.} $\abs{A \cap B} = \abs A + \abs B - \abs{A \cup
B} \geq \abs A + \abs B - n > \tfrac{4n}3 - n = \tfrac n3$,
则
\[
\abs{A \cap B \cap C} \geq \abs{A \cap B} + \abs C - n >
\tfrac n3 + \tfrac{2n}3 - n = 0 .
\]

\textbf{18.} 注意 $m > 2n/3$。对 $T_1, T_2, T_3 \subseteq S$
且基数为 $m$ 者，由第 17 问，存在点 $x \in T_1
\cap T_2 \cap T_3$；于是 $x \in \operatorname{conv}(T_i)$
对每个 $i$ 成立：$\mathcal F$ 中任意三个成员相交。
该族为有限族（$S$ 的子集仅有有限多个），且由
\omterm{def:b2:affine:convex}{凸集}构成，故由赫利定理（第 12 问）得 $c \in
\bigcap_{\abs T = m}\operatorname{conv}(T)$。

\textbf{19.} 设某闭半平面 $H \ni c$
含少于 $n/3$ 个 $S$ 中的点。其补集 $U$
为开半平面，是\omterm{def:b2:affine:convex}{凸集}，且 $\abs{S \cap U} > 2n/3$，
从而 $\abs{S \cap U} \geq \floor{2n/3} + 1 = m$；取 $T
\subseteq S \cap U$ 使 $\abs T = m$。于是
$\operatorname{conv}(T) \subseteq U$，由 $U$ 为\omterm{def:b2:affine:convex}{凸集}，故
$c \in \operatorname{conv}(T) \subseteq U$，与
$c \in H$ 矛盾。因此包含 $c$ 的每个闭半平面
至少含 $n/3$ 个点：$c$ 为中心点。

\textbf{20.} 取边长为 $L$ 的正三角形，并取
$\varepsilon = L/100$。设 $c$ 为任一点；我们给出一个
包含 $c$ 且至多含 $k$ 个点的闭半平面。

\emph{情形 1：$c$ 到某一顶点的距离不超过 $L/10$，不妨设该顶点为 $B$。}
自 $c$ 到 $A$、$C$ 的方向，与
$B \to A$、$B \to C$ 的方向相比，偏差至多
$\arcsin\bigl(\tfrac{L/10}{9L/10}\bigr) = \arcsin\tfrac19$,
于是两者所成夹角 $\leq \tfrac\pi3 + 2\arcsin\tfrac19 <
\tfrac{2\pi}3$。\emph{情形 2：$c$ 到各顶点的距离
均 $> L/10$。}若 $c$ 在三角形内，
则 $u_A, u_B, u_C$ 为 $c$ 指向各顶点的方向，其间三个角隙
之和为 $2\pi$，故某角隙 $\leq 2\pi/3$；
若 $c$ 在三角形外，三个方向落在某一开半平面
内，其中两者的夹角 $<
\pi/2$。总之必有两方向（不妨指向 $X$ 与 $Y$）
夹角 $\leq 2\pi/3$；记 $w$ 为其单位角平分向量，
于是 $\langle u_X, w\rangle, \langle u_Y, w\rangle \geq
\cos\tfrac\pi3 = \tfrac12$。对圆盘
中任一点 $b$（圆盘以 $X$ 为心）：
\[
\langle \vect{cb}, w\rangle \geq
\tfrac12\norm{\vect{cX}} - \varepsilon > 0,
\]
由于 $\norm{\vect{cX}} \geq L/10 > 2\varepsilon$（对
$Y$ 亦然），开半平面 $\{\langle \vect{cx},
w\rangle > 0\}$ 将两簇点都包含在内。其
闭补集含 $c$，且至多含第三簇中的
$k$ 个点。故平面上任何点都不能保证多于 $n/3$：
结合第 19 问，中心点常数恰为 $1/3$。

\textbf{21.} 将三个角排序；最大者 $\theta$，
满足 $\theta \geq \pi/3$（三角之和为 $\pi$）。若
$\theta \geq \pi/2$，设在 $R$ 处，令 $M$ 为
对边 $\intcc PQ$ 的中点。由中线公式（$\vect{RM}
= \tfrac12(\vect{RP} + \vect{RQ})$，展开，
并借助余弦定理消去 $\langle\vect{RP}, \vect{RQ}\rangle$）
得
\[
\norm{\vect{RM}}^2 = \tfrac12\norm{\vect{RP}}^2 +
\tfrac12\norm{\vect{RQ}}^2 - \tfrac14\norm{\vect{PQ}}^2 \leq
\tfrac12\norm{\vect{PQ}}^2 - \tfrac14\norm{\vect{PQ}}^2 =
\tfrac14\norm{\vect{PQ}}^2,
\]
利用 $\norm{\vect{PQ}}^2 = \norm{\vect{RP}}^2 +
\norm{\vect{RQ}}^2 - 2\langle\vect{RP}, \vect{RQ}\rangle \geq
\norm{\vect{RP}}^2 + \norm{\vect{RQ}}^2$（内积
为 $\leq 0$）。故直径为 $\intcc PQ$ 的圆盘，半径
$\leq \tfrac12 < \tfrac1{\sqrt3}$，包含全部三点
（退化的共线三点归入 $\theta = \pi$ 的情形）。若
$\theta < \pi/2$，三角形为锐角三角形；由正弦定理
外接圆半径为 $R_c = a/(2\sin\theta)$，其中 $a \leq 1$
为 $\theta$ 所对的边，且 $\theta \in
\intco{\pi/3}{\pi/2}$ 得 $\sin\theta \geq \sqrt3/2$，故
$R_c \leq 1/\sqrt3$，外接圆盘即可。

\textbf{22.} 对 $p \in S$，令 $K_p = \overline D(p,
1/\sqrt3)$，为紧\omterm{def:b2:affine:convex}{凸集}。设 $p, q, s$ 为 $S$ 中任意三点，
两两距离 $\leq 1$，由第 21 问可得一
半径为 $1/\sqrt3$ 的圆盘包含它们：其圆心位于
$K_p \cap K_q \cap K_s$ 中。由紧形式的赫利定理（第 14 问，
允许任意族），存在 $c \in \bigcap_{p\in
S}K_p$：每个 $p \in S$ 都在半径为 $1/\sqrt3$、中心为 $c$ 的圆盘内，即 $S
\subseteq \overline D(c, 1/\sqrt3)$。此即
平面上的荣格定理。

\textbf{23.} 设 $G$ 为\omterm{def:b2:affine:barycenter}{重心}，$\sum_i\vect{GA_i} = 0$，
故
\[
\sum_i\norm{\vect{OA_i}}^2 = \sum_i\norm{\vect{OG} +
\vect{GA_i}}^2 = 3\norm{\vect{OG}}^2 + 2\Bigl\langle
\vect{OG}, \sum_i\vect{GA_i}\Bigr\rangle +
\sum_i\norm{\vect{GA_i}}^2,
\]
中间项为零，此即莱布尼茨恒等式。对
边长为 $1$ 的正三角形，$\norm{\vect{GA_i}} =
1/\sqrt3$（即高 $\sqrt3/2$ 的三分之二），故
$\sum_i\norm{\vect{GA_i}}^2 = 1$。若 $\overline D(O, r)$
包含三顶点，则 $3r^2 \geq
\sum_i\norm{\vect{OA_i}}^2 = 3\norm{\vect{OG}}^2 + 1 \geq 1$：
$r \geq 1/\sqrt3$，等号成立时必有 $O = G$，且
三条距离皆等于 $r$——即外接圆盘。荣格
常数 $1/\sqrt3$ 是最优的。

\textbf{24.} \emph{$\R^n$ 中的赫利定理：若 $C_1, \dots, C_m$
（$m \geq n + 1$）为 $\R^n$ 的凸子集，且其中任意 $n + 1$
个相交，则全体相交。} 基础情形 $m = n +
2$：取 $x_i \in \bigcap_{j\neq i}C_j$；拉东引理
（\cref{exo:b2:affine:7}）将 $x_1, \dots, x_{n+2}$ 分为
两块 $I \mid J$，并有公共\omterm{def:b2:affine:convex}{凸包}点 $z$，且对每个
$k$，不含 $k$ 的那一块由
$C_k$ 中的点组成，故由凸性得 $z \in C_k$，与第
11 问相同。归纳步：对 $m \geq n + 2$，将 $C_m,
C_{m+1}$ 替换为 $C_m \cap C_{m+1}$；新族中任一 $(n+1)$ 元组，
若含该交集成员，则相当于旧族中的 $n + 2$
个集合，由基础情形处理；其余元组
由题设覆盖。由归纳假设即
得结论。

\textbf{25.} 线性代数只在此处出现一次：$n + 2$ 个向量位于
$(n+1)$ 维的数对空间（总权重，
加权位置）中必线性相关——即
仿射相关性。其系数的符号将诸点分为
两个拉东块，并将一个线性关系化为
两个非负\omterm{def:b2:affine:barycenter}{重心}相等的等式。凸性恰被用到
两次：一是赫利归纳步中（$C_k$ 中各点的\omterm{def:b2:affine:convex}{凸包}
仍在 $C_k$ 中），二是应用中（半平面与
圆盘皆为凸）。平面的维数仅在
输入拉东引理的点数 $4 = 2 + 2$ 中出现，即
赫利题设中的“$3 = 2 + 1$”；其余一切
均与维数无关，正如第 24 问所证实。在 $\R^n$ 中
各常数变为：赫利数为 $n + 1$；中心点常数
$\frac1{n+1}$（每个有限集皆有一点，使得过该点的每个闭
半空间都至少含有其 $\geq
\frac1{n+1}$ 的点）；荣格半径
$\sqrt{\frac{n}{2(n+1)}}$，对直径为 $1$ 的集合——等于
$1/\sqrt3$，当 $n = 2$ 时即为此值。
\end{solution}
